\documentclass[9pt,a4paper,twocolumn,twoside]{tau-class/tau}
\usepackage[spanish,es-nodecimaldot,es-noindentfirst]{babel}

\doctype{Reporte Ejecutivo Integrado}
\title{Reporte de performance y competencia: Psicotel con Terapias}

\author[a,1]{Equipo de Analitica Oderbiz}
\affil[a]{Unidad de Inteligencia Comercial}

\dates{Compilado el \today}
\docinfo{Documento generado desde datos internos en formato \texttt{llm\_context\_report.page.v1} e investigacion competitiva publica. Toda inferencia esta etiquetada como \texttt{inferencia}.}
\footinfo{Reporte Oderbiz}
\theday{\today}
\organization{Oderbiz}
\leadauthor{Equipo Oderbiz}

\begin{abstract}
Entre agosto de 2024 y abril de 2026, la cuenta analizada alcanzo 832{,}006 impresiones y 644{,}214 de alcance con \$781.91 de inversion acumulada (\texttt{observado}). El embudo evidencia fuerte volumen en clics unicos (8{,}835) y conversaciones iniciadas (614), pero baja salida efectiva a enlaces externos (97 outbound clicks), con costo por outbound de \$8.06 (\texttt{observado}). En terminos competitivos, el mercado local de Loja muestra oferta fragmentada de psicologos individuales, centros con portafolio amplio y listados agregadores con fuerte comparabilidad de precio/servicios (\texttt{observado}). La principal oportunidad es escalar creatividad y segmentacion orientadas a mensajes de alta intencion y derivacion controlada a WhatsApp/sitio para convertir demanda en consultas cerradas (\texttt{inferencia}). Se propone un plan P1/P2/P3 enfocado en eficiencia de embudo, control geografico y mejora de conversion conversacional.
\end{abstract}

\keywords{Meta Ads, psicologia infantil, Loja, performance marketing, inteligencia competitiva}

\begin{document}
\maketitle
\thispagestyle{firststyle}

\section{Contexto del negocio y periodo analizado}

Analisis de la pagina \texttt{1506380769434870}, cuenta \texttt{act\_131112367482947}, con ventana historica \texttt{maximum} (2024-08-13 a 2026-04-21) (\texttt{observado}).

El inventario de campanas muestra predominio de objetivo \texttt{MESSAGES}, complementado por \texttt{OUTCOME\_ENGAGEMENT}, con presencia puntual de \texttt{OUTCOME\_TRAFFIC}/\texttt{LINK\_CLICKS} (\texttt{observado}).

\section{Baseline interno}

\subsection{KPIs principales}
\begin{itemize}
  \item Inversion total: \$781.91; impresiones: 832{,}006; alcance: 644{,}214 (\texttt{observado}).
  \item CTR global: 1.58\%; CPM global: \$0.94 (\texttt{observado}).
  \item Funnel: 8{,}835 clics unicos, 614 conversaciones iniciadas, 541 primeras respuestas (\texttt{observado}).
  \item Calidad de trafico: 97 outbound clicks, \$8.06 por outbound, \$0.09 por clic unico (\texttt{observado}).
\end{itemize}

\subsection{Tendencia temporal}
Se identifican tres fases de comportamiento:
\begin{enumerate}
  \item \textbf{Fase de eficiencia alta} (Q4-2024 a Q1-2025): CPM y CPC consistentemente bajos, con CTR sostenido alrededor de 1.1\%--2.8\% (\texttt{observado}).
  \item \textbf{Fase de estabilidad media} (Q2-2025 a Q4-2025): CTR relativamente estable con oscilaciones por variacion creativa y frecuencia de pauta (\texttt{inferencia}).
  \item \textbf{Fase de presion de costos} (Q1-2026 a Q2-2026): se incrementa CPM (hasta rango 1.2--1.7) y CPC en varios dias, sugiriendo mayor competencia de subasta o fatiga de audiencias (\texttt{inferencia}).
\end{enumerate}

\subsection{Embudo y calidad}
La relacion entre 8{,}835 clics unicos y 97 outbound clicks evidencia que la arquitectura actual optimiza interaccion inicial, pero con baja transferencia a propiedad digital externa (\texttt{observado}). Para un modelo centrado en mensajeria, esto puede ser funcional; no obstante, limita atribucion posterior y escalamiento de conversion medible fuera de la plataforma (\texttt{inferencia}).

\subsection{Sena geo y demografica}
\textbf{Geo}: Loja concentra mayor volumen y resultados de mensajeria (575 conexiones de mensajeria), seguida por Zamora Chinchipe y contribuciones menores en otras provincias (\texttt{observado}).

\textbf{Edad}: 25--44 aporta el grueso de volumen y resultados de conversacion; 45+ muestra mejor costo por clic en algunos tramos pero con menor profundidad de mensajeria en relacion al volumen (\texttt{observado}+\texttt{inferencia}).

\section{Panorama competitivo}

\subsection{Lectura de mercado local}
La competencia en Loja combina tres frentes:
\begin{itemize}
  \item \textbf{Profesionales individuales} con posicionamiento por especialidad (psicologia clinica, infantil, TCC, lenguaje) (\texttt{observado}).
  \item \textbf{Centros con cartera integral} (psicologia clinica, educativa, neuropsicologia, psiquiatria) y precios publicados por paquete (\texttt{observado}).
  \item \textbf{Marketplaces/directorios} que facilitan comparacion por ciudad, modalidad y precio, elevando sensibilidad a reputacion y claridad de oferta (\texttt{observado}).
\end{itemize}

\subsection{Implicaciones competitivas}
\begin{itemize}
  \item Mensajes de valor deben diferenciarse por resultado concreto (hitos de desarrollo, plan terapeutico, seguimiento familiar), no solo por disponibilidad (\texttt{inferencia}).
  \item Publicar propuesta y ruta de atencion clara reduce fuga hacia listados comparativos (\texttt{inferencia}).
  \item Integrar evidencia social local (casos, testimonios verificados, contenidos educativos) aumenta defensa frente a competidores de alto volumen en directorios (\texttt{inferencia}).
\end{itemize}

\section{Diagnostico integrado}

\subsection{Fortalezas}
\begin{itemize}
  \item Cobertura robusta con bajo CPM historico para categoria de servicios (\texttt{observado}).
  \item Capacidad probada para generar conversaciones (614 iniciadas, 541 primeras respuestas) (\texttt{observado}).
  \item Mix de campanas activo y continuidad operativa de pauta (\texttt{observado}).
\end{itemize}

\subsection{Brechas}
\begin{itemize}
  \item Muy baja derivacion outbound vs volumen de clic unico (\texttt{observado}).
  \item Nomenclatura de campanas inconsistente (\textit{campana sin nombre}) dificulta aprendizaje y control analitico (\texttt{observado}).
  \item Menor trazabilidad de calidad por cohorte creativa/segmento (\texttt{inferencia}).
\end{itemize}

\subsection{Riesgos}
\begin{itemize}
  \item Incremento de costos de subasta y caida de eficiencia en periodos recientes (\texttt{observado}+\texttt{inferencia}).
  \item Dependencia fuerte de mensajeria sin capa robusta de medicion CRM/lead lifecycle (\texttt{observado}+\texttt{inferencia}).
  \item Exposicion a comparacion de precio/servicio en directorios locales (\texttt{observado}).
\end{itemize}

\subsection{Oportunidades}
\begin{itemize}
  \item Escalar segmentos 25--44 con creatividades por problema especifico (lenguaje, conducta, aprendizaje) (\texttt{inferencia}).
  \item Construir embudo mixto: mensaje + derivacion cualificada a WhatsApp/sitio con seguimiento (\texttt{inferencia}).
  \item Explotar dominancia geografica en Loja y replicar clusters adyacentes de alto potencial (\texttt{inferencia}).
\end{itemize}

\section{Plan de accion priorizado}

\subsection{P1 (alto impacto / bajo esfuerzo)}
\begin{itemize}
  \item Estandarizar naming de campanas, conjuntos y anuncios (objetivo--audiencia--oferta--fecha).
  \item Duplicar mejores anuncios de mensajes con 3 variantes de copy orientadas a motivo de consulta.
  \item Implementar guion de respuesta inicial para reducir tiempo de primera respuesta y elevar tasa de cita.
  \item Crear KPI semanal minimo: costo por conversacion iniciada, costo por primera respuesta, costo por cita agendada.
\end{itemize}

\subsection{P2 (alto impacto / esfuerzo medio-alto)}
\begin{itemize}
  \item Diseñar embudo dual: \texttt{MESSAGES} + \texttt{TRAFFIC} con pagina de aterrizaje enfocada en conversion.
  \item Construir taxonomia de audiencias: padres por rango de edad del hijo + intereses de desarrollo infantil.
  \item Lanzar experimento creativo por format (testimonio, educacion, caso-problema-solucion) con lectura quincenal.
  \item Incorporar medicion de calidad post-mensaje (contacto efectivo, asistencia, continuidad de terapia).
\end{itemize}

\subsection{P3 (experimental)}
\begin{itemize}
  \item Modelos lookalike basados en conversaciones de alta calidad (\texttt{inferencia}).
  \item Colaboraciones de contenido con entidades locales (educacion/salud) para autoridad de marca (\texttt{inferencia}).
  \item Biblioteca de micro-contenido diagnostico preventivo para captacion organica + paid amplification (\texttt{inferencia}).
\end{itemize}

\section{Fuentes y limitaciones}

\subsection{Fuentes internas usadas}
\begin{itemize}
  \item \texttt{report\_metadata}
  \item \texttt{page\_overview}
  \item \texttt{timeseries\_daily}
  \item \texttt{funnel}
  \item \texttt{traffic\_quality}
  \item \texttt{geo}
  \item \texttt{demographics}
  \item \texttt{campaigns\_available}
\end{itemize}

\subsection{Fuentes competitivas publicas}
\begin{itemize}
  \item \url{https://psicologiaymente.com/psicologos/2069561/dra-marizol-jimenez-luzon}
  \item \url{https://psicologiaymente.com/psicologos/2062410/nicole-loaiza}
  \item \url{https://www.psico.org/ec/loja/terapia-de-pareja}
  \item \url{https://www.psico.org/ec/directorio/loja}
  \item \url{https://www.terappio.com/ec/family-therapist/ecuador/loja}
  \item \url{https://tecnologicosudamericano.edu.ec/bienestar-estudiantil/centro-psicoterapeutico-la-cabana/}
\end{itemize}

\subsection{Limitaciones detectadas}
\begin{itemize}
  \item No se ejecuto \texttt{apify-market-research} por ausencia de evidencia operativa de token/creditos en esta corrida.
  \item No se dispuso de CRM cerrado (citas atendidas, ticket, recurrencia), por lo que el analisis de revenue downstream queda parcial.
  \item No hay campo de competencia directa validada por el cliente; se uso proxy de presencia publica local (\texttt{inferencia}).
\end{itemize}

\end{document}
